\documentclass[12pt]{article}
\usepackage[utf8]{inputenc}
\usepackage[T1]{fontenc}
\usepackage[margin=1in]{geometry}
\usepackage{setspace}
\usepackage{xcolor}
\usepackage[colorlinks=true, urlcolor=blue!60!black, linkcolor=blue!60!black]{hyperref}
\setlength{\parindent}{0em}
\setlength{\parskip}{0.6em}
\renewcommand{\baselinestretch}{1.0}
\begin{document}
\begin{center}
{\LARGE \textit{Shelby County v. Holder}, 570 U.S. 529 (2013)\par}
\vspace{0.5em}
{\large Case Brief --- Two-Page Edition\par}
\end{center}

\section{Memory Jogger}
The Voting Rights Act's \S~4(b) coverage formula, based on decades-old data, can no longer subject States to \S~5 preclearance (\href{https://www.law.cornell.edu/supremecourt/text/12-96}{Cornell LII}).

\section{Detailed Case Facts}
Section 5 of the 1965 Act required covered jurisdictions to obtain federal preclearance of any voting change; \S~4(b) defined coverage by 1964-era tests and low turnout. \textit{Katzenbach} (1966) upheld the scheme. In 2006, after a 15,000-page record, Congress extended it 25 years without updating the formula. \textit{Northwest Austin} (2009) had warned that current burdens must be justified by current needs.

\section{Procedural History}
Shelby County, Alabama, sued Attorney General Eric Holder for a declaration that \S\S~4(b) and 5 were facially unconstitutional. The district court and a divided D.C. Circuit upheld the law; the Supreme Court reversed in 2013.

\section{Judicial Votes}
5--4. Roberts wrote for Scalia, Kennedy, Thomas, and Alito; Thomas concurred; Ginsburg dissented, joined by Breyer, Sotomayor, and Kagan.

\section{Holding}
Section 4(b) is unconstitutional because it rests on 40-year-old data and eradicated practices. The Court did not rule on \S~5 and left \S~2 intact, but preclearance became inoperable absent a new formula.

\section{Analysis of Opinions}
\textbf{Majority:} The ``fundamental principle of equal sovereignty'' and the States' power over elections made preclearance ``strong medicine.'' ``[N]early 50 years later, things have changed dramatically'': registration gaps had closed and objections had nearly vanished. A formula must be relevant to current conditions, and Congress did not use its record to write a new one. \textbf{Thomas:} \S~5 itself is unconstitutional. \textbf{Ginsburg:} Congress has broad enforcement power; discrimination evolved into ``second-generation barriers''; and the majority ignored the record and silently overruled \textit{Katzenbach} on equal sovereignty.

\section{Examples: Future Applications}
\textbf{Same-Side:} a formula covering States that used poll taxes before 1964 is invalid. \textbf{Opposite-Side:} a court may order ``bail-in'' preclearance under \S~3(c) after finding intentional discrimination. \textbf{Fence-Sitter:} whether equal sovereignty limits a sports-wagering law that favors two States.

\section{Critique}
``Equal sovereignty'' is not in the text, and the Reconstruction Amendments altered federalism. The Court also left open whether \textit{Katzenbach} or \textit{Boerne} supplies the standard, which \textit{Louisiana v. Callais} (2026) resolved in favor of \textit{Boerne}.

\section{Key Quotations}
\begin{quote}
``Throwing out preclearance when it has worked and is continuing to work to stop discriminatory changes is like throwing away your umbrella in a rainstorm because you are not getting wet.''
\end{quote}
\noindent --- Justice Ginsburg, dissenting, 570 U.S. at 590.

\end{document}
